\setchapterpreamble[u]{\margintoc}
\chapter{Measurement of the Transfer Function}
\labch{tf-measurement}

If we have an optical system and want to model the way light passes through that system, either in order to enhance our
    understanding of the structure (source) that causes the signal or for any other reason, we need to determine one of
    the transfer functions introduced in the previous part. % TODO: reference <14-11-23> 
Of the transfer functions, it is most intuitive to directly measure the \ac{PSF} or more specifically the
    magnitude of the \ac{PSF}.
We disregard the phase change that the optical system may introduce, because in most set-ups it is actually impossible 
    to directly measure the imaginary part of the \ac{PSF}\sidenote[][*+1]{It is possible in some cases to determine the 
    imaginary part (or the phase) from the magnitude algorithmically which will be discussed in 
    \refsec{phase-retrieval}.}, since the light sensors\sidenote[][*+5]{In our case the \ac{CMOS} sensor.} only measure 
    the energy that hits the sensor elements which is identical for all phases of the \ac{EM} wave since it only depends 
    on the intensity of the \ac{EM} wave (\refchshort{light}). % TODO: Add closer reference. <14-11-23> 


\section{Image Formation Model}%
\labsec{image-formation-model}

% TODO: First step, noise. Measure noise from the background of the beads image. <14-11-23> 
First off, let us devise an image formation model based on the way, we expect the light from the source to manifest 
    itself in the resulting signal that is captured by the sensor.


\subsection{Illumination}%
\labsubsec{illumination}
To determine the transfer function of the optical system, means to measure it for a uniformly illuminated focal plane.
However, we do not have data directly from a regular uniformly illuminated wide field system.
% TODO: Reference specific section <22-11-23> 
Our images are acquired with the harmonic illumination (\refchshort{SR-microscopy}). 
This means that for a still\sidenote[][*-9]{Meaning that the source of the light is stationary in contrast to \emph{in 
    vivo} biological acquisitions, where this is not the case.} image, we have 9 wide field images with 3 different 
    pattern orientations and 3 different phases\sidenote[][*-6]{In accordance with the method that is used to generate 
    the patterns in the case of our \ac{SIM} setup, we are assuming that the phases differ by \(2\pi/3\) from one 
    another and do not determine them independently.} for each of those.

% TODO: Change to margin figure? The differences are not so big anyway <22-11-23> 
\begin{marginfigure}[*-2]
    \marginpng{beads-9-patterns-by-orientations-and-phases}
    \caption{Acquisitions of 9 wide field images with harmonic patterned illuminations. The orientations are labelled
    1, 2 and 3 and the phases \(-\), \(0\) and \(+\).}
    \labfig{beads-LR-9-images}
\end{marginfigure}

In order to account for this discrepancy, we will imitate the uniform illumination by summing all the images acquired
    with different illuminations together (this will also reduce noise (\refsubsecshort{noise})).
Due to additivity of the optical system, this results in the same signal, as if an acquisition was made with the sum
    of the illumination patterns.
For three patterns in one orientation \(I_-\), \(I_0\) and \(I_+\), where the phases are \(\phi_0 - 2\pi/3\), \(\phi_0\)
    and \(\phi_0 + 2\pi/3\), let us, without loss of generality, assume that \(\phi_0 \equiv 0\), since we can move the 
    origin as we please, simplifying the sum \(I_\Sigma\) along the line of the wave vector of the harmonic pattern
    i.e. \(\vec{r} = \alpha(\vec{k}_I)^{\transpose}\), gives

    % TODO: Graphic of the line in the pattern <22-11-23> 

    \begin{equation*}
        \begin{split}
            I_\Sigma (\vec{r}) &= \sum_{i=1}^3 1 - \frac{m_i}{2} \cos\ab(2\pi \vec{k}_I \cdot \vec{r} 
                + \frac{2\pi i}{3}) \\
                               &= 3 - \sum_{i=1}^3 \frac{m}{2} \cos\ab(2\pi \vec{k}_I \cdot \vec{r} + 
                                \frac{2\pi i}{3}) \\
                               &\begin{split}
                                   = 3 &-\frac{m}{2} \cos\ab(2\pi \vec{k}_I \cdot \vec{r})
                                    \underbrace{\sum_{i=1}^3\ab(\cos\frac{2\pi i}{3})}_{=0} \\
                                       &- \frac{m}{2} \sin\ab(2\pi  \vec{k}_I \cdot \vec{r})
                                    \underbrace{\sum_{i=1}^3\ab(\sin\frac{2\pi i}{3})}_{=0},
                                \end{split}
        \end{split}
    \end{equation*}
    where, in the second step, we made the assumption that \(m_1 = m_2 = m_3\)\sidenote[][*+1]{This assumption holds 
    precisely if the illumination pattern is generated by a diffraction grating, which is our case, and also 
    theoretically if laser interference is used.}.

\begin{marginfigure}[*-14]
\marginpdftex{sum-of-patterns-z-x-axes-for-one-orientation}
\caption{Sum of the illumination pattern intensity along the wave vector for the first orientation, results in a uniform 
    intensity.}
\labfig{sum-of-ips}
\end{marginfigure}

% TODO: Show that along another line, this only projects the wave vector on the line, and the phases do not change so
% the sum is also uniform <22-11-23> 

% TODO: Section about the images <16-11-23> 
% TODO: Figure: Add reconstructed (HR) image figure and LR image figures <16-11-23> 

% TODO: Second step, illumination... The sum <16-11-23> 
    % TODO: Side figure of the z-x axis cut of the illumination intensity cut <16-11-23> 
    % TODO: Note about additivity and why the 3× illuminated signal is the same as if we would have illuminated it all at
    % once <16-11-23> 

\subsection{Noise}%
\labsubsec{noise}

As a first step towards developing a faithful image formation model, we need to consider noise.

The mode of noise, that we will be particularly interested in is \emph{additive noise}, for which the strength of the 
    noise is added to the noise free signal:
\begin{equation} \labeq{additive-noise-model}
    \sensed(\vec{x}) = \overbrace{\noisefree(\vec{x})}^{\text{signal}} + \underbrace{\noise(\vec{x})}_{\text{noise}}.
\end{equation}
In \refeq{additive-noise-model} the noise \(\noise(\vec{x})\) has a specified distribution\sidenote{It is not 
    uncommon to choose to model the noise by a Gaussian distribution (\(\distribution{D} \equiv \distribution{N}\)) even
    if the Gaussian distribution is physically unrealistic for the given noise source. This simplification is often made 
    for the purpose of simplifying computation.}
\begin{equation*}
    \noise(\vec{x}) \sim \distribution{D}.
\end{equation*}
An important property to consider when characterising noise from different sources is \emph{signal independence}, i.e. 
    whether the strength of the noise is independent from the intensity of the signal
\begin{equation*}
    \left\{ \noise(\vec{x}) \mid \noisefree(\vec{x})  \right\} = \noise(\vec{x}).
\end{equation*}

There are many sources of noise that should be considered in the acquisition model.
The most prominent of these sources are caused by the image sensor (\ac{CMOS})\sidenote{Especially in low \ac{SNR}
    acquisitions.}. %  and the fluorescent proteins (\ac{EGFP}).
It is also common to regard the variations in the ambient (background) illumination i.e. the light that is not part of 
    the source we are measuring as noise.
The sources, that are the greatest contributors to the degradation of the expected signal are:
\begin{description}
    % TODO: Add figure of the relation of the SNR ratio against the signal strength in this source of noise <24-11-23> 
    % TODO: Add figure with the Poisson distribution <25-11-23> 
    \item[Photon noise]
        Inherent to the quantized nature of energy transmission of light. 
        Sensors measure irradiance by counting the number of discrete photons that land on the sensor in a given time
            interval. 
        The arrivals of individual photons are independent and can therefore be regarded as a Poisson process.
        The number \(N\) of photons detected in the time interval \(\Delta\tau\) can thus be described by the
            distribution
            \begin{equation*}
                \distribution{P}\left[ N = k \right] = \frac{e^{-\lambda\Delta\tau}(\lambda\Delta\tau)^k}{k!},
            \end{equation*}
            where the rate parameter \(\lambda\) corresponds to the expected incident photon count per unit time, i.e.
            the expected number of photon detections per a unit of time.
        It is dependent on the source luminance.
        Since \(\E N = \Var(N) = \lambda\Delta\tau\), the ratio of signal to noise increases with
            \(\sqrt{\lambda\Delta\tau}\) and so the corruption due to this source of noise is greatest in situations 
            with low signal levels.
        For larger counts of photons, we can use the central limit theorem to approximate the photon noise by a Gaussian
            additive noise \sidecite{hasinoff2021}, whose variance and mean both depend on the signal level
            \begin{gather*}
                \sensed(\vec{x}) = \noisefree(\vec{x}) + \noise(\vec{x}) \\
                \noise(\vec{x}) \sim \distribution{N}\left(\lambda(\vec{x})\cdot\Delta\tau, 
                                                         \lambda(\vec{x})\cdot\Delta\tau \right).
            \end{gather*}
    \item[Dark current]
        Semiconductor detectors in the image sensor are susceptible to thermal agitation\sidecite{aguet2009}, which can 
            cause the photocapacitors to charge even in the absence of light (hence the name dark current) leading to 
            false signals.
        Due to the phenomenon of dark current, the sensors are often cooled to reduce this effect.
        Noise from this source is signal independent and additive.
    \item[Background noise]
        Additive and independent noise from the ambient light that reaches the sensor or unwanted light caused by other
            effects\sidenote{An example can be the lagged emission of the fluorescent proteins.}.
        This effect of ambient light from outside sources is most prominent for low energy light, i.e. light with low 
            frequencies (e.g. infrared light), which is more prone to scattering and harder to effectively shade out.
        In the case of biological sensing however, a majority of what we consider background noise from the light that 
            is far from the focal plane and defocus causes it to blur and appear as ambient light with locally (meaning 
            in a small bounded region) limited variation.
        This is the principal reason for the use of a \ac{TIRFM} setup.
    \item[Auxiliary sources]
        A variety of causes mainly in the imaging sensor that are produced at the read-out time\sidenote{The phases
            after the photo activation of the sensor elements (consisting of amplification,
            digitization of the analog signal etc.).}.
        Notably the charge transfer in the \ac{CMOS} sensor can be the cause noise.
        Auxiliary sources can be either signal dependent or independent.
\end{description}

If we know or assume a distribution of additive noise, it is possible to estimate its parameters using parts of the
    image with approximately or theoretically constant intensity (e.g. background regions).
Importantly, it is often the case that we assume, or in approximation assume, that the distribution of the additive 
    noise is Gaussian and we can estimate the mean and the standard deviation of the noise using the \ac{MLE}
\begin{gather*}
    A \coloneq \left\{ \vec{x} \mid \noisefree(\vec{x}) \approx c  \right\} \\
    \hat{\mu}_{\noise} = \frac{1}{|A|} \sum_{\vec{x} \in A} (\sensed(\vec{x}) - c) \\
    \hat{\sigma^2}_{\noise} = \frac{1}{|A|} \sum_{\vec{x} \in  A} (\sensed(\vec{x}) - \hat{\mu}_{\noise})^2.
\end{gather*}
Moreover, if the limitation to the theorized noise distribution is to be avoided, the distribution can be eyeballed from
    the histogram or the histogram, kernel or other empirical estimate of the distribution of the intensities in \(A\)
    can be used.

\begin{marginfigure}
    % FIX: Font sizes on figures should be adjusted and this figures has to have a different layout (GridLayout) <21-11-23> 
    % \marginpng{beads-leakage-on-otsu-threshold}
    \png{beads-leakage-on-otsu-threshold}
    \caption{
        The Otsu threshold leads to lower intensities of some beads to be passed as background. 
        On the left, side, Otsu threshold (\(\approx~0.13\)) is used and on the right, the it is lowered by \(0.07\) to 
            \(\approx~0.06\). 
        The top image shows a larger area of the difference of these two.
    }
    \labfig{beads-leakage-on-otsu-threshold}
\end{marginfigure}

In our case, the premise that allows us to select such a set \(A\) is the lack of other signal sources besides the
    beads.
The problem therefore reduces to the distinction of the beads from the background, for which an Otsu
    threshold\sidecite[*-6]{otsu1979} can be used, however it is necessary to lower the threshold to reduce the lower
    intensities (\reffigshort{beads-leakage-on-otsu-threshold}), that are due to the beads fluorescence.

\begin{marginfigure}
    \marginpng{noise-distribution-hist-from-background}
    \caption{Histogram of the background noise. It is apparent that the threshold cuts off some of the data. This is due
    to a low \ac{SNR}.}
    \labfig{noise-distribution-hist-from-background}
\end{marginfigure}
The histogram of the background data that we consider to be the distribution of \emph{background} noise is in
    \reffig{noise-distribution-hist-from-background}.

% TODO: MLE noise parameter estimates with a truncated distribution. That is due to the threshold, the full distribution
% is not visible and therefore the full MLE is not totally correct. <22-11-23> 

If we consider the background noise as the true noise distribution, disregarding other sources of noise, we can 
    estimate the \ac{SNR}\sidenote{
    A popular measure used in signal processing to quantify the relative power of noise. For an image it can 
        be defined as
        \begin{equation}
            \labeq{snr}
        \SNR=\frac{\sum_{\substack{u \in \hat{M} \\ v \in \hat{N}}}\ab|\noisefree(u,v)|^2}{\sum_{\substack{u \in \hat{M} \\ v \in
    \hat{N}}}\ab|\noise(u,v)|^2}
        \end{equation}
} of the image as well.


% TODO: Add SNR estimation and reason, why the noise is reduced if we add multiple acquisitions together <25-11-23> 

It is practical to assume that the, yet unknown, source signal \(\source(\vec{x})\), is directly proportional to the
    concentration of the fluorescent molecules, which in turn is directly proportional to their volume in focus.
For this, it is necessary to disregard the noise sources, which arise at the source level, but this is a reasonable
    proposition, owing to the fact, that most of the noise can be attributed to the sensor and is added to the sensed 
    signal after the pass through the optical apparatus and the background, which can be effectively removed\sidenote{It 
    is important to note that only the power of the noise, linked to the mean value, is effectively reduced, not the 
    variation as such. Methods of removing the background noise generally remove the low spatial frequencies, however 
    the background noise is also prone to the photon shot phenomenon and other signal dependent noise sources, which 
    reside even after the background removal.}.

\subsection{The Transfer Function}%
\labsubsec{the-transfer-function}

Next step is to consider how the source signal \(\source(\vec{x})\) is affected by the pass through the optical system.
This has been developed in \refch{transfer-functions}.
As a reminder, in general, for a linear system \(H\), the Fredholm equation of the first kind applies \(\forall \vec{r}
    \in \R^2\)
    \begin{multline} \labeq{non-aplanatic-psf-light-formation}
        \noisefree(\vec{r}) = H\ab[\source(\vec{r})] = H\ab[\iint_{\R^2} \source(\vec{x}) \dirac(\vec{r} - \vec{x}) \odif{\vec{x}}] 
            = \iint_{\R^2} \source(\vec{x}) H\ab[\dirac(\vec{r} - \vec{x})] \odif{\vec{x}}  \\
            = \iint_{\R^2} \source(\vec{x}) \psf(\vec{r} -\vec{x}; \vec{r}) \odif{\vec{x}},
    \end{multline}
    or without the focal plane simplification \(\forall\vec{r} \in \R^3\)
    \begin{equation*}
        \noisefree(\vec{r}) = \iiint_{\R^3} \source(\vec{x}) h_3(\vec{r} - \vec{x}; \vec{r}) \odif{\vec{x}},
    \end{equation*}
    where we are using the same notation as in \refeq{additive-noise-model} and
    \begin{equation*} \labeq{PSF-in-fredholm}
        \psf(\vec{r} - \vec{x}; \vec{r}) \coloneq H\ab[\dirac(\vec{r} - \vec{x})]
    \end{equation*}
    denotes the \ac{PSF}.
Measuring the \ac{PSF} \(h\) that is location dependent like in \refeq{non-aplanatic-psf-light-formation} is badly
    conditioned, so it is necessary to simplify the transfer.
For this, the assumption that the system is aplanatic\sidenote{An optical system, or more specifically the lens
    % TODO: Closer reference <24-11-23> 
    configuration is aplanatic, if it is free of both spherical and coma aberrations (\refchshort{transfer-functions}).} 
    is made.
This assumption makes the \ac{PSF} focal plane location independent i.e. \(\forall\vec{x},\vec{r} \in \R^2\)
    \begin{equation*} \labeq{aplanar-psf-light-formation}
        \psf(\vec{r} - \vec{x}; \vec{r}) = \psf(\vec{r} - \vec{x})
    \end{equation*}
     in the same \(z\)-offset plane from the focal plane.
However it is not true that under the assumption of an aplanatic system that \ac{PSF} \(h_3\) is identical as a 
    function \(\R^3 \to \R\) for the whole imaged volume.
\marginnote{
    The aplanatic system assumption is not fully justified and there is a slight location dependence in practice.
    This issue can be addressed by estimating the \ac{PSF} on patches of the sensed image separately.
}
The lack of spherical aberration and coma only fixes \(h\) on the \(x\) and \(y\) axes, so for 
    \(\vec{x},\vec{r} \in \R^3\)
    \begin{equation*}
       h_3(\vec{r} - \vec{x}; \vec{r})  \equiv h_3(\vec{r} - \vec{x}; z).
    \end{equation*}


For the location invariant \(h\), we can rewrite \refeq{non-aplanatic-psf-light-formation} as
    \begin{equation}
        \labeq{convolution-psf}
        \noisefree(\vec{r}) = \iint_{R^2} \source(\vec{x}) \psf(\vec{r} - \vec{x}) \odif{\vec{x}} = 
            (\source * \psf)(\vec{r}),
    \end{equation}
    where \(*\) denotes the convolution operation.
The locality\sidenote{Locality is understood here as the property that if \(\supp(\psf) \subset 
    \ball_\alpha\), then for a source \(\source\) for which \(\supp(\source) \subset \ball_\beta\) we
    there exists a \(\gamma\) such that \(\supp(\source * \psf) \subset \ball_\gamma.\)} of the convolution is key to making the 
    \ac{PSF} more intuitive and easier in practice to measure.
It is, because unlike the \ac{OTF}, we are not forced to make the estimate from the full image or strenuously map the
    frequency domain of the parts used for measurement onto the frequencies of the acquisition image.
% TODO:  What changes when the PSF changes...  <22-11-23, kunzaatko> 
% TODO: What changes when the OTF changes <22-11-23> 
Consequently, parts that are easiest to work with can be selected.
% TODO: Another condition for this to work is that the PSF has a bounded support with a bound that is limited by the
% region sizes that we select. This should be argued. It is not the case but the PSF has is small enough to consider
% this as true <27-11-23> 
This can be beneficial in eliminating noisy or inconsistent data, or data, that our model does not consider.


\subsection{The Source}
\labsubsec{the-source}
In order to devise \(\psf\) from \refeq{convolution-psf}, it is necessary to determine \(\source(\vec{r})\), for the
    parts that are used for the measurement.
Consider what it would would mean for the acquisition, if \(\source= \sum_{\vec{a} \in A} \dirac_{\vec{a}}\).
The sensed image, where noise is disregarded, would be
    \begin{equation}
        \labeq{noisefree-with-sum-of-dirac-sources-psf-general}
        \noisefree(\vec{r}) = \ab(\source * \psf)(\vec{r}) = \ab(\ab(\sum_{\vec{a} \in A} \dirac_{\vec{a}}) * 
            \psf)(\vec{r}) = \sum_{\vec{a} \in A} \psf(\vec{r} - \vec{a}),
    \end{equation}
    and if adapted to the more natural notation, where we emphasize, which point underlines the term of the expansion in
    the image formation formula, i.e. \(\psf(\vec{r} - \vec{a}) \eqqcolon \psf_{\vec{a}}
    (\vec{r})\), \refeq{noisefree-with-sum-of-dirac-sources-psf-general} changes to
    \begin{equation}
        \labeq{noisefree-with-sum-of-dirac-sources-psf-natural}
        \noisefree(\vec{r}) = \sum_{\vec{a} \in A} \psf_{\vec{a}}(\vec{r}).
    \end{equation}
\marginnote[*-5]{\refeq{noisefree-with-sum-of-dirac-sources-psf-natural} is also applicable for the non-aplanatic 
    transfer, but it is important to realize that in that case the index \(\vec{a}\) of the \ac{PSF} does not denote 
    only a change in coordinates (translation by \(\vec{a}\)), but also the change in the function \(\psf_{\vec{a}}\) 
    itself.} %  This can be used alongside partitioning the sensed image into regions and measuring the relation of the 
    % \ac{PSF} to the location with in the \refsec{location-dependent-psf}}

In practice, this type of source is generated by imaging microspheres most commonly made from plastics, stained with 
    fluorescent dyes.
This is a source that is often used in practice by biologists for many purposes of calibrating equipment, because it is 
    possible to manufacture with a high degree of precision and modeling it is straightforward to model.
An important aspect that is essential to be able to measure the \ac{PSF} in a direct manner is that the diameters of the
    beads is sub-diffraction-resolution\sidenote[][*-3]{It is possible, and even beneficial if the diameter is not
    sub-pixel-resolution, but it does not make a difference in the principle of the \ac{PSF} measurement.}.


If the diameter were not to be sub-diffraction-resolution the source signal would be different from the
    \(\dirac_{\vec{a}}\) of \refeq{noisefree-with-sum-of-dirac-sources-psf-general}, but rather some other support 
    limited source for which a model \(\hat{\source}_{\vec{a}}\) would be necessary to determine, and the noise free
    signal \(\noisefree\) would not be direct copies of the \ac{PSF} but
    \begin{equation}
        \labeq{non-sub-diffraction-resolution-source}
        \noisefree(\vec{r}) = \sum_{\vec{a} \in A} \ab(\hat{\source}_{\vec{a}} * \psf_{\vec{a}})(\vec{r})
    \end{equation}
    instead.
Deriving the \ac{PSF} \(\psf\) from the signal in \refeq{non-sub-diffraction-resolution-source} would require 
    deconvolution of the individual components \((\hat{\source}_{\vec{a}} * \psf_{\vec{a}})\) with the assumed model 
    \(\hat{\source}_{\vec{a}}\) and due to noise, this would necessarily mean that we would need to make further 
    assumptions either about the deconvolved \(\hat{\psf}_{\vec{a}}\) or the noise.

The acquisition that is used for the \ac{PSF} estimation is done by imaging a suspension of \emph{TetraSpeck\TM 
    microspheres, \(\qty{0.1}{\micro\metre}\)}, where the diameter is 
    \(\qty{0.099 \pm 0.008}{\micro\metre}\)\sidecite{smith2020}, which is safely below the diffraction resolution.

% TODO: Figure with selected beads marked from the acquisition <28-11-23> 

% TODO: Either in this section or in the section about the Abbe resolution limit, there should be the two point
% resolution definition <27-11-23> 


\section{\ac{PSF} by Averaging Fluorescent Beads}%
\labsec{psf_by_averaging_fluorescent_beads}

It is common in literature to select one bead from the centre of the acquisition image and pronounce it as
    the \ac{PSF} of the optical system.
% TODO: References (these references should be those which are used elsewhere 
% (perhaps in the SIM reconstruction papers)) <28-11-23> 
There is ample room for improvement in this approach. 
Noteworthy drawbacks of this estimate are:
    \begin{itemize}
        \item noise contained in a single bead sample can be substantial,
        \item possibility of incorrect centring,
        \item disregard of the relation on the \(z\)-offset from the focal plane of the particular bead 
            sample used,
        \item negligence towards the potential dependence of the sampled bead position.
    \end{itemize}

    from multiple beads (\refsubsecshort{noise}).

% TODO: Motivational parts for averaging <28-11-23> 

% TODO: Make this through algorithm pseudocode typography <28-11-23> 
\begin{enumerate}
    \item Select centres of isolated beads with the highest maximum intensity
    \item Cut out an array patch around the bead
    \item Remove background % TODO: What am I doing for background removal <28-11-23> 
    \item Find sub-pixel the centres of the beads by \ac{LSE} fitting a Gaussian to the sample array
    \item Shift by sub-pixel parts of the beads centres to fix the bead centres to an integer pixel value via Fourier
        shift theorem
    \item Average the patches to obtain the \ac{PSF} estimation
\end{enumerate}
